% Part: turing-machines
% Chapter: undecidability
% Section: introduction

\documentclass[../../../include/open-logic-section]{subfiles}

\begin{document}

\olfileid{tur}{und}{int}
\olsection{مقدمة}

قد يُظن أن استحالة حساب بعض الدوال أمر ظاهر، ولو اقتصرنا على
الدوال الحسابية؛ لكثرتها وما يعرض في سلوكها من التعقيد. فمنها ما
يُعرّف باضمحلال الجسيمات المشعة، أو بغيره من الظواهر الفوضوية أو
العشوائية. ولبيان ذلك، نفرض لحظة نبتدئ منها عد فترات طول كل منها
ثانية، ونجعل $f(n)$ عدد جسيمات الكون التي تضمحل في الفترة
رقم~$n$ بعد تلك اللحظة. فهذه دالة يبدو أن حسابها مما لا يُطمع فيه.

غير أن خفاء طريق الحساب علينا ليس برهانًا على امتناعه. فقد تكون
الدالة، على شدة تعقيدها في الظاهر، قابلة للحساب بالمعنى المجرد.
ولنفرض، على سبيل المثال، أن للكون عمرًا منتهيًا، وأنه ينكمش في مستقبل
بعيد إلى نقطة واحدة، كما تتنبأ بعض النظريات الكونية. فلا يكون بعد
اللحظة التي فرضناها إلا عدد منتهٍ من الثواني تكون فيها $f(n)$
معرّفة، وإن عظم هذا العدد غاية العظم. وكل دالة يقتصر مجال تعريفها
على مدخلات منتهية العدد قابلة للحساب؛ إذ يمكن جمع مخرجاتها في جدول
واحد ضخم، أو تضمينها مخطط انتقال حالات لآلة تورنغ بالغة الكبر.

ومن الدوال ما يخصنا بالنظر، وهو ما كانت قيمه أجوبة عن أسئلة لا
تحتمل إلا نعم أو لا. فالسؤال «أعدد $n$ أولي؟» تقابله الدالة
\[
\fn{isprime}(n) = \begin{cases}
  1 & \text{إذا كان $n$ أوليًا}\\
  0 & \text{خلاف ذلك.}
  \end{cases}
\]
ونسمي سؤال نعم أو لا \emph{قابلًا للتقرير فعليًا} متى كانت دالته
المقابلة، ذات القيمتين $1/0$، قابلة للحساب فعليًا.

وإثبات وجود دوال لا تُحسب فعليًا، أو مسائل لا تُقرر فعليًا، يقتضي
أولًا تعيين نموذج للحساب، ثم إقامة البرهان على عجزه عن بعض الدوال
أو المسائل. فنبين مثلًا أن آلات تورنغ لا تحسب كل دالة ولا تقرر كل
مسألة. فإذا ثبت ذلك، جاز لنا الاستناد إلى أطروحة تشرتش--تورنغ:
فليس العجز حينئذ خاصًا بهذه الآلات، بل يلزم أن لا يكون لأي إجراء
فعلي قدرة على حساب كل دالة.

وأصل الحجة في هذه النتائج السلبية أن آلات تورنغ نفسها تقبل إسناد
الأعداد إليها. وأيسر سبيل إلى ذلك تعدادها: نعيّن طريقة لكتابة الآلات
وبرامجها، ثم نسرد تلك الكتابات على نظام معلوم. ومتى تحقق إمكان هذا
التعداد، كفى النظر في الأعداد الأصلية لإثبات وجود دوال لا تحسبها آلة
تورنغ. فإن مجموعة الدوال من $ \Nat$ إلى~$\Nat$، بل مجموعة الدوال من
$\Nat$ إلى $\{0, 1\}$ وحدها، !!{nonenumerable}؛ وأما مجموعة الدوال
القابلة للحساب بآلة تورنغ فهي !!{denumerable}، إذ نستطيع تعداد
الآلات جميعًا.

ولا يقتصر الأمر على إثبات الوجود؛ بل نستطيع تعيين دوال لا تقبل
الحساب ومسائل لا تقبل التقرير، وإقامة البرهان على كل منها. ومن ذلك
\emph{مسألة التوقف}. فآلة تورنغ لها وصف منتهٍ هو قائمة تعليماتها،
وهذا الوصف، أي برنامج الآلة، يصلح دخلًا لآلة أخرى. ومن ثم يصح أن
نطلب إلى آلة تورنغ تقرير أسئلة عن سائر آلات تورنغ. وأهمها هنا:
«أتتوقف الآلة المعطاة في نهاية المطاف إذا ابتدأت بالدخل~$n$؟»
ولو وجدت آلة تقرر ذلك لأشبهت جهازًا لضبط الجودة، يضمن ألا تقع الآلات
في حلقات لا تنتهي ونحوها. لكن تورنغ أثبت أن وجود هذه الآلة ممتنع:
فلا آلة واحدة تأخذ وصف آلة تورنغ~$M$ وعددًا~$n$ دخلًا، ثم تتوقف
في جميع الأحوال بخرج $1$ أو $0$، بحسب توقف $M$ أو عدم توقفها إذا
ابتدأت بالدخل~$n$.

ومتى ثبت أن مسألة بعينها لا تقبل التقرير، أمكن اتخاذها سبيلًا إلى
إثبات امتناع تقرير مسائل أخرى. ومن أشهرها السؤال:
«أصيغة الرتبة الأولى~$!A$ صادقة منطقيًا؟» فلا توجد آلة تورنغ تأخذ
صيغة رتبة أولى~$!A$ وتضمن التوقف بخرج $1$ أو $0$ بحسب كون
$!A$ صادقة منطقيًا أو غير صادقة كذلك. وقد عُرف طلب إجراء فعلي لحل
هذا السؤال في تاريخ المنطق باسم «مسألة القرار» بإطلاق. ولهذا يقال:
مسألة القرار غير قابلة للحل. وأثبت تورنغ وتشرتش هذه النتيجة
مستقلين، في وقتين متقاربين؛ فسميت أيضًا مبرهنة تشرتش--تورنغ.

\end{document}
